\section{The Abel--Ruffini contradiction}
\label{sec:abel-proof}

We now combine the three ingredients:

\begin{enumerate}
\item a polynomial split by a radical tower has solvable Galois group
  (Corollary~\ref{cor:split-rad-solvable});
\item the generic quintic has Galois group $S_5$
  (Theorem~\ref{thm:generic-galois-S5});
\item $S_5$ is not solvable
  (Theorem~\ref{thm:S5-not-solvable}).
\end{enumerate}

\subsection{No radical tower splits the generic quintic}

\begin{theorem}[Generic quintic obstruction]
\label{thm:generic-quintic-obstruction}
Let $K,L,P$ be as in Section~\ref{sec:generic-quintic}.  No radical tower over
$K$ contains all five roots $x_1,\ldots,x_5$ of $P$.
\end{theorem}

\begin{proof}
Suppose a radical tower contained all roots.  Since $L=K(x_1,\ldots,x_5)$,
the tower would contain the splitting field $L$.  By
Corollary~\ref{cor:split-rad-solvable}, $\Gal(L/K)$ would be solvable.  By
Theorem~\ref{thm:generic-galois-S5}, this group is isomorphic to $S_5$, which
is not solvable by Theorem~\ref{thm:S5-not-solvable}.  Contradiction.
\end{proof}

The proof is intuitionistically valid.  The assumption of a radical tower is
used to construct an explicit finite solvability certificate; the $S_5$
calculation refutes that certificate.  No double-negation elimination is used.

\subsection{From a universal formula to a radical tower}

Suppose there were a universal formula by radicals for monic quintics
\[
 T^5+a_1T^4+a_2T^3+a_3T^2+a_4T+a_5.
\]
The formula may contain finitely many branches, field operations, and radical
instructions.  Specialize its coefficient inputs to
\[
 a_1=-e_1,
 \quad a_2=e_2,
 \quad a_3=-e_3,
 \quad a_4=e_4,
 \quad a_5=-e_5
\]
in the field $K$.  Follow the branch guards, and adjoin the finitely many
radical witnesses in program order.  The resulting run is a radical tower.
Completeness of the formula supplies all five roots and therefore a complete
factorization of $P$.  This contradicts
Theorem~\ref{thm:generic-quintic-obstruction}.

This specialization argument covers formulas with exceptional cases.  A
formula that is truly universal must have a valid branch for the generic
coefficient tuple.  It also covers formulas whose output roots are unordered,
whose branches repeat roots, or whose radicals have several possible values:
any successful run still generates a radical tower containing a complete root
list.

\begin{theorem}[Abel--Ruffini in degree five]
There is no universal complete formula, built from rational constants, field
operations, finite equality case distinctions, and extraction of finitely
many radicals, that solves every degree-five equation over characteristic-zero
fields.
\end{theorem}

\subsection{Every degree greater than five}

For $n\ge5$, put
\[
 P_n(T)=T^{n-5}P(T).
\]
This is monic of exact degree $n$.  Its root multiset contains
$x_1,\ldots,x_5$ together with $n-5$ zero roots.  If a radical tower split
$P_n$, it would contain all $x_i$ and hence split $P$, contradicting
Theorem~\ref{thm:generic-quintic-obstruction}.

\begin{corollary}[Abel--Ruffini in every degree at least five]
For each $n\ge5$, no universal complete radical formula exists for all
polynomials of degree $n$.
\end{corollary}

The padding argument deliberately proves a statement about complete solution.
The padded polynomial has the easy root zero; producing that one root does not
solve the equation.

\subsection{Why the coefficient field is legitimate}

A reader may ask whether a formula intended for complex coefficients must be
tested over the rational-function field $K$.  There are two equivalent
answers.

\paragraph{Formal-expression answer.}
A universal radical formula is an identity in the coefficient indeterminates.
Its operations are algebraic and therefore make sense over every
characteristic-zero field where its denominators are nonzero and its guarded
branches are followed.  The generic field is precisely the universal test
object for such identities.

\paragraph{Specialization answer.}
If one insists on complex coefficients, embed a finitely generated
specialization of the generic data into a sufficiently large characteristic-
zero field and preserve the finitely many nonzero denominators.  A successful
universal formula would then yield the same radical-tower consequence.  An
alternative route, already present in ProveIt, uses explicit rational
polynomials with nonsolvable Galois group.  Section~\ref{sec:explicit-quintic}
records one elementary certificate route.

\subsection{Strong and weak formulations}

The following implications are useful when comparing formal statements.

\[
\begin{array}{c}
\text{a fixed coefficient formula returns all roots}\\[0.2em]
\Downarrow\\[-0.1em]
\text{each input polynomial has a complete radical expression}\\[0.2em]
\Downarrow\\[-0.1em]
\text{the generic quintic has a complete radical factorization}\\[0.2em]
\Downarrow\\[-0.1em]
\Gal(P/K)\text{ is solvable}.
\end{array}
\]

The proof refutes the bottom statement.  Therefore it refutes every stronger
notion above it, including nonuniform per-root expressions when they are
required for all five roots.  This is why the argument is insensitive to the
surface syntax chosen for ``formula by radicals.''
